% Source: https://artofproblemsolving.com/community/c3h2256267_peace09s_1365th_post
% Stable identity: aops-topic-2256267-q02
Both sides have degree $\dfrac{n(n+1)}{2}=1+2+\cdots+n$, but they are not always equal.

\textbf{It holds for $n=1$ and $n=2$:} $1+x=x+1$, and $(x+1)(x^2+1)=x^3+x^2+x+1$.

\textbf{It fails for every $n\ge3$.} Substitute $x=1$. The left side has $\dfrac{n(n+1)}{2}+1$ terms, each equal to $1$, so it equals $\dfrac{n(n+1)}{2}+1$; the right side equals $2^n$. For $n=3$ these are $7$ and $8$. For $n\ge3$ we always have $2^n>\dfrac{n(n+1)}{2}+1$: this is true for $n=3$, and if it holds for $n$, then
\[2^{n+1}>n(n+1)+2\ge\dfrac{(n+1)(n+2)}{2}+1,\]
since $n(n+1)+2-\left(\dfrac{(n+1)(n+2)}{2}+1\right)=\dfrac{(n+1)(n-2)}{2}+1>0$. So the two polynomials take different values at $x=1$ and are not equal.

(Another way to see it: expanding the product, the coefficient of $x^k$ counts the subsets of $\{1,2,\dots,n\}$ with sum $k$; for $n\ge3$ the coefficient of $x^3$ is $2$, from $\{3\}$ and $\{1,2\}$, while on the left it is $1$.)

So the statement is \textbf{false}; the identity holds exactly for $\boxed{n=1\text{ and }n=2}$.
